\chapter{Required Template for the Design of Each Project's Module}

\newcommand{\taskA}{\textit{Task A}}
\newcommand{\taskB}{\textit{Task B}}
\newcommand{\taskC}{\textit{Task C}}

\section*{Notes}

\textbf{The final design document delivered for each module must not exceed 2 pages per student task and 7 pages in total.}

\textbf{The examples given in each section/subsection are not complete/exhaustive!}

% ============================= SECTION =============================
\section{Module Overview and Requirements}

This section will contain the requirements for each module. Read it carefully and don't change it or remove it.

% ============================= SECTION =============================
\section{Design Details}

% ------------------------- SUB-SECTION -------------------------
\subsection{Design Description}

Here you should describe the design of your solution. This should be the main and most substantial part of the document.

The solution should be described in terms of the main data structures and the functions or interfaces you will introduce or modify.

\begin{itemize}
    \item the key invariants (properties that must always be true) of your code and how they maintained;
    \item how shared state is protected and concurrent operations are coordinated;
    \item how and where resources are allocated, used, and released;
    \item how errors and edge cases are handled.
\end{itemize}

Recommendations:
\begin{itemize}
    \item Use pseudocode to describe algorithms
    \item Use diagrams to illustrate the architecture your solution
    \item Don't list every minor helper function
    \item Don't include existing declarations without explanation
    \item Don't provide a line-by-line description of the implementation
\end{itemize}

\begin{tcolorbox}[title={Example}]
\paragraph{Timer} We store a list of active timers, each with a list of waiting threads. The list of timers is checked periodically to see if any of them expired. When a timer expires, all waiting threads are released.

To block the threads waiting for a timer, we use executive events, which allow blocking threads with \lstinline{ExEventWaitForSignal}.
To unblock all waiting threads, we initialize the event as \lstinline{ExEventTypeNotification} and call \lstinline{ExEventSignal}.

The events are initialized and uninitialized together with their timer, with the observation that they should be signaled before the timer is destroyed.

\paragraph{Priority Scheduler. Fixed Priority Scheduler} We change how threads are inserted into the waiting lists, so they are ordered descending by priority,
using \lstinline{InsertOrderedList}. To ensure that the highest priority threads are running at any given time,
we will send an IPI (interprocessor interrupt) to invoke the scheduler on each CPU.
Specifically, we need to do this when:
\begin{itemize}
    \item a thread is inserted into the global ready list, when a new thread is created or when a thread is unblocked from a timer, mutex or event.
    \item the priority of threads changes, due to priority modifications or priority donation.
\end{itemize}

\paragraph{Priority Scheduler. Priority Donation} Priority donation will be structured like a tree with two kinds of nodes:
\begin{itemize}
    \item thread - each thread has a list of held mutexes
    \item mutex - each mutex has a list of waiting threads
\end{itemize}
The mutex nodes only play an ``intermediary'' role, passing the priority of the threads waiting for the respective mutex
to the holder of the mutex. This donation can happen recursively, with a chain like
$Thread_A \xrightarrow{Waits} Mutex_1 \xrightarrow{Donates} Thread_B \xrightarrow{Waits} Mutex_2 \xrightarrow{Donates} Thread_C$,
which must be handled by ``propagating'' the priority down to the running thread (in this case $Thread_C$).

We also need to consider the cases when a thread holds multiple mutexes, choosing the highest priority of all donations, and when a mutex has multiple
waiting threads, again, choosing the highest priority.
\end{tcolorbox}

\subsubsection{\taskA}

\studentTodo{Complete the text for \taskA}

\subsubsection{\taskB}

\studentTodo{Complete the text for \taskB}

\subsubsection{\taskC}

\studentTodo{Complete the text for \taskC}

% ------------------------- SUB-SECTION -------------------------
\subsection{Interactions Between Components (Student Tasks)}

Here you should describe each case where the code/logic for a given component (student task) depends on other components (student tasks).
After you identify the interactions, we recommend defining an interface for each component, so the users
of the given component (i.e. your teammates who implement the other tasks) don't need to know its internal details.

Given that the implementation is in C, an interface will be a header file containing function headers and type definitions.
Before starting the actual implementation work (i.e. modifying \texttt{.c} files), each team member should create or extend one
or more (existing) header (\texttt{.h}) files, that expose the functions and types that will be used by their teammates
when implementing the other modules.

\begin{tcolorbox}[title={Example}]
\paragraph{Timer} There is no direct interaction between the Timer and the Scheduler or Donation. There is an indirect interaction,
because the timer makes use of executive events, which need to take into consideration thread priorities when releasing the waiting threads.

\paragraph{Priority Scheduler and Priority Donation}
The priority donation component depends on the scheduler, because it must be notified whenever thread priorities change.

The priority scheduler should expose a function to the priority donation component, so the priority donation
can ``notify'' the scheduler that a priority changed during donation and it may need to reschedule threads.
\end{tcolorbox}

\studentTodo{Complete the text as a team, defining a subsubsection for each interaction, or to note the lack of interactions}

% ------------------------- SUB-SECTION -------------------------
\subsection{Design Decisions, Motivations and Alternatives}

Here you should explain the key design decisions you made and the motivations behind them.
Discuss any alternative approaches you considered and why you rejected them.

Recommendations:
\begin{itemize}
    \item discuss realistic alternatives, that you actually considered
    \item describe each alternative briefly, without going into significant details
\end{itemize}

\begin{tcolorbox}[title={Example}]
\paragraph{Timer} The timers are checked in the system timer interrupt (\lstinline{ExSystemTimerTick}), which is simple and already provided, but
its main disadvantage is that it's not precise. An alternative that should have significantly better precision is to use ACPI or PIT timer, by
programming the next interrupt to be exactly at the trigger time of the next timer. This requires both a significant amount of
research to understand how these mechanisms and their implementation works, so we discarded this option.

\paragraph{Priority Scheduler. Fixed Priority Scheduler} The data structure used to store waiting/blocked threads was implemented as a doubly-linked circular list, and uses \lstinline{InsertOrderedList},
because this required the least amount of implementation work, and is easy to understand.
The main alternative is using a more efficient data structure, for example, an array of doubly-link circular lists, with one list for each priority level. While this would be more efficient, it would also use more memory
and would take more effort to implement.

\paragraph{Priority Scheduler. Priority Donation} In the ``donation tree'', we always propagate the priority of the donor thread to the root, because
it's easy to prove correct. We considered if this is really needed, and if we could stop early when we see a mutex or thread with a higher priority
down the chain, but discarded the idea because we weren't sure that we covered all edge cases.
\end{tcolorbox}

\subsubsection{\taskA}

\studentTodo{Complete the text for \taskA}

\subsubsection{\taskB}

\studentTodo{Complete the text for \taskB}

\subsubsection{\taskC}

\studentTodo{Complete the text for \taskC}


% ============================= SECTION =============================
\section{Implementation Considerations}

Here you should describe the practical aspects that you needed to take into consideration during the implementation.
This can include:

\begin{itemize}
    \item main challenges and issues encountered during the implementation
    \item interesting bugs encountered, how you found the root cause, how you fixed them and how you could prevent similar bugs in the future
    \item test cases that required more effort to pass than you expected
\end{itemize}

\begin{tcolorbox}[title={Example}]
\paragraph{Timer} We noticed that all timer tests already passed before starting the implementation, but they started to fail once
we introduced partial modifications into the code.

\paragraph{Priority Scheduler. Fixed Priority Scheduler}
To force other CPUs to yield, we used \lstinline{SmpIpiSendToAllIncludingSelf} and initially called \lstinline{ThreadYield} inside the IPI,
but we noticed that the system would deadlock after a while. We then understood that functions that invoke the scheduler can't be
called inside an interrupt. So to fix the issue, we instead set the \lstinline{YieldOnInterruptReturn} field of the current CPU.
\end{tcolorbox}

\subsubsection{\taskA}

\studentTodo{Complete the text for \taskA}

\subsubsection{\taskB}

\studentTodo{Complete the text for \taskB}

\subsubsection{\taskC}

\studentTodo{Complete the text for \taskC}